% =========================================================
% Rational Density and Order Facts
% =========================================================

\section{Rational Density and Order Facts}
\label{sec:rational-density-and-order}

The rational numbers are not merely numerous; they are tightly packed.  Between
any two distinct rationals one can find another, and small rational
perturbations can always be made without leaving $\mathbb{Q}$.

% ---------------------------------------------------------
% Theorem: Between Any Two Rational Numbers Lies Another Rational
% Label: thm:between-any-two-rationals-is-a-rational
% ---------------------------------------------------------
\begin{tcolorbox}[colback=thmbox, colframe=thmborder, arc=2pt,
  left=6pt, right=6pt, top=4pt, bottom=4pt,
  title={\small\textbf{Theorem (Between Any Two Rational Numbers Lies Another Rational)}},
  fonttitle=\small\bfseries]
\begin{theorem}[Between Any Two Rational Numbers Lies Another Rational]
\label{thm:between-any-two-rationals-is-a-rational}
If $r,s\in\mathbb{Q}$ and $r<s$, then there exists $t\in\mathbb{Q}$ such that
\[
r<t<s.
\]
\end{theorem}
\end{tcolorbox}

\begin{remark*}[Interpretation]
This is the order-density of $\mathbb{Q}$ in itself.  Rational numbers have no
adjacent neighbors.
\end{remark*}

\begin{remark*}[Dependencies]
Basic arithmetic closure of $\mathbb{Q}$ under addition and division by $2$.
\end{remark*}

\begin{remark*}[Proof idea]
Take the midpoint $t=(r+s)/2$.  It is rational and lies strictly between $r$
and $s$.
\end{remark*}

\begin{remark*}[Proof]
\hyperref[prf:between-any-two-rationals-is-a-rational]{Go to proof.}
\end{remark*}

% ---------------------------------------------------------
% Corollary: Q Is Dense in Itself
% Label: cor:q-dense-in-itself
% ---------------------------------------------------------
\begin{tcolorbox}[colback=thmbox, colframe=thmborder, arc=2pt,
  left=6pt, right=6pt, top=4pt, bottom=4pt,
  title={\small\textbf{Corollary ($\mathbb{Q}$ Is Dense in Itself)}},
  fonttitle=\small\bfseries]
\begin{corollary}[$\mathbb{Q}$ Is Dense in Itself]
\label{cor:q-dense-in-itself}
The subset $\mathbb{Q}$ is order dense in itself.
\end{corollary}
\end{tcolorbox}

\begin{remark*}[Dependencies]
\hyperref[thm:between-any-two-rationals-is-a-rational]{Theorem~\ref*{thm:between-any-two-rationals-is-a-rational}}.
\end{remark*}

\begin{remark*}[Proof idea]
This is exactly the content of the theorem above, expressed as a structural
property of the ordered set $\mathbb{Q}$.
\end{remark*}

\begin{remark*}[Proof]
\hyperref[prf:q-dense-in-itself]{Go to proof.}
\end{remark*}

% ---------------------------------------------------------
% Corollary: The Rational Numbers Have No Adjacent Points
% Label: cor:q-has-no-adjacent-points
% ---------------------------------------------------------
\begin{tcolorbox}[colback=thmbox, colframe=thmborder, arc=2pt,
  left=6pt, right=6pt, top=4pt, bottom=4pt,
  title={\small\textbf{Corollary (The Rational Numbers Have No Adjacent Points)}},
  fonttitle=\small\bfseries]
\begin{corollary}[The Rational Numbers Have No Adjacent Points]
\label{cor:q-has-no-adjacent-points}
There do not exist rational numbers $r<s$ such that no rational number lies
strictly between them.
\end{corollary}
\end{tcolorbox}

\begin{remark*}[Dependencies]
\hyperref[thm:between-any-two-rationals-is-a-rational]{Theorem~\ref*{thm:between-any-two-rationals-is-a-rational}}.
\end{remark*}

\begin{remark*}[Proof idea]
Apply the theorem to the interval $(r,s)$.
\end{remark*}

\begin{remark*}[Proof]
\hyperref[prf:q-has-no-adjacent-points]{Go to proof.}
\end{remark*}

% ---------------------------------------------------------
% Theorem: Q Has No Minimum Positive Element
% Label: thm:q-has-no-minimum-positive-element
% ---------------------------------------------------------
\begin{tcolorbox}[colback=thmbox, colframe=thmborder, arc=2pt,
  left=6pt, right=6pt, top=4pt, bottom=4pt,
  title={\small\textbf{Theorem ($\mathbb{Q}$ Has No Minimum Positive Element)}},
  fonttitle=\small\bfseries]
\begin{theorem}[$\mathbb{Q}$ Has No Minimum Positive Element]
\label{thm:q-has-no-minimum-positive-element}
For every $q\in\mathbb{Q}$ with $q>0$, there exists $r\in\mathbb{Q}$ such that
\[
0<r<q.
\]
\end{theorem}
\end{tcolorbox}

\begin{remark*}[Interpretation]
Positive rationals can always be halved.  There is no smallest positive
rational number.
\end{remark*}

\begin{remark*}[Dependencies]
Basic arithmetic closure of $\mathbb{Q}$ under division by $2$.
\end{remark*}

\begin{remark*}[Proof idea]
Take $r=q/2$.
\end{remark*}

\begin{remark*}[Proof]
\hyperref[prf:q-has-no-minimum-positive-element]{Go to proof.}
\end{remark*}

% ---------------------------------------------------------
% Lemma: Upper/Lower Perturbation Lemma (Rational Version)
% Label: lem:rational-upper-lower-perturbation
% ---------------------------------------------------------
\begin{tcolorbox}[colback=thmbox, colframe=thmborder, arc=2pt,
  left=6pt, right=6pt, top=4pt, bottom=4pt,
  title={\small\textbf{Lemma (Upper/Lower Perturbation Lemma, Rational Version)}},
  fonttitle=\small\bfseries]
\begin{lemma}[Upper/Lower Perturbation Lemma, Rational Version]
\label{lem:rational-upper-lower-perturbation}
If $q\in\mathbb{Q}$ and $\varepsilon>0$, then there exist $r,s\in\mathbb{Q}$
such that
\[
q-\varepsilon < r < q < s < q+\varepsilon.
\]
\end{lemma}
\end{tcolorbox}

\begin{remark*}[Dependencies]
\hyperref[thm:q-has-no-minimum-positive-element]{Theorem~\ref*{thm:q-has-no-minimum-positive-element}},
basic closure of $\mathbb{Q}$ under addition and subtraction.
\end{remark*}

\begin{remark*}[Proof idea]
Choose a small positive rational $\delta<\varepsilon$ and set
$r=q-\delta$, $s=q+\delta$.
\end{remark*}

\begin{remark*}[Proof]
\hyperref[prf:rational-upper-lower-perturbation]{Go to proof.}
\end{remark*}

% ---------------------------------------------------------
% Lemma: Rational Perturbation Lemma
% Label: lem:rational-perturbation-lemma
% ---------------------------------------------------------
\begin{tcolorbox}[colback=thmbox, colframe=thmborder, arc=2pt,
  left=6pt, right=6pt, top=4pt, bottom=4pt,
  title={\small\textbf{Lemma (Rational Perturbation Lemma)}},
  fonttitle=\small\bfseries]
\begin{lemma}[Rational Perturbation Lemma]
\label{lem:rational-perturbation-lemma}
If $q\in\mathbb{Q}$ and $\varepsilon>0$, then there exists $r\in\mathbb{Q}$
such that
\[
0<|r-q|<\varepsilon.
\]
\end{lemma}
\end{tcolorbox}

\begin{remark*}[Interpretation]
Every rational point can be moved a little without leaving $\mathbb{Q}$.
\end{remark*}

\begin{remark*}[Dependencies]
\hyperref[lem:rational-upper-lower-perturbation]{Lemma~\ref*{lem:rational-upper-lower-perturbation}}.
\end{remark*}

\begin{remark*}[Proof idea]
Take either the upper or lower perturbation from the preceding lemma.
\end{remark*}

\begin{remark*}[Proof]
\hyperref[prf:rational-perturbation-lemma]{Go to proof.}
\end{remark*}
